\textbf{月考复习}

By Deralive

\textbf{一、单选题}

1．函数 $f(x) = x^{2} - \text{cos}x$ 在区间 $\left\lbrack 0,\text{π} \right\rbrack$ 上的平均变化率为（~~~）

\begin{quote}
    A． $- \text{π} - \frac{2}{\text{π}}$  B． $- \text{π}$  C． $\text{π}$
    D． $\text{π} + \frac{2}{\text{π}}$
\end{quote}

2．已知函数  $f(x) = \frac{1}{3}x^{3} - x，$ 则
$\lim_{\bigtriangleup x \rightarrow 0}\frac{f(3) - f(3 - 2 \bigtriangleup x)}{\bigtriangleup x} =$ （~~~）

\begin{quote}
    A．2 B．4 C．8 D．16
\end{quote}

3．已知函数 $f(x)$ 在 $x = x_{0}$ 处可导，且 $\lim_{\Delta x \rightarrow 0}\frac{f\left( x_{0} - \text{Δ}x \right) - f\left( x_{0} \right)}{\text{Δ}x} = 2025$ ，则 $f^{\prime}\left( x_{0} \right) =$ （~~~~）

\begin{quote}
    A． $- 2025$  B．0 C．1 D．2025
\end{quote}

4．设\emph{f}(\emph{x})是可导函数，若 $\lim_{\Delta x \rightarrow 0}\frac{f(3 - 3\Delta x) - f(3)}{\Delta x} = 3$ ，则 $f^{\prime}(3)$ （~~~~）

\begin{quote}
    A． $- 1$  B． $- \frac{1}{3}$  C． $\frac{1}{3}$  D．1
\end{quote}

5．某物体沿直线运动，位移 $y$ （单位： $\text{m}$ ）与时间 $t$ （单位： $\text{s}$ ）的关系为 $y(t) = 1 - t + t^{2}$ ，则该物体在 $t = 3\text{ }\text{s}$ 时的瞬时速度是（~~~~）

\begin{quote}
    A． $2\text{ }\text{m}/\text{s}$  B． $3\text{ }\text{m}/\text{s}$
    C． $4\text{ }\text{m}/\text{s}$  D． $5\text{ }\text{m}/\text{s}$
\end{quote}

6．如图，已知函数 $f(x)$ 的图象在点 $P\left( 2,f(2) \right)$ 处的切线为 $l$ ，则 $f(2) + f^{\prime}(2) =$ （~~~~）

\begin{quote}
    A． $- 3$  B． $- 2$  C．1 D．2
\end{quote}

7．函数 $f(x)$ 的图象如图所示， $f^{\prime}(x)$ 为函数 $f(x)$ 的导函数，下列数值排序正确的是（~~~）

\begin{quote}
    A． $0 < f^{\prime}(2) < f^{\prime}(3) < f(3) - f(2)$
    B． $0 < f^{\prime}(3) < f(3) - f(2) < f^{\prime}(2)$

    C． $0 < f^{\prime}(3) < f^{\prime}(2) < f(3) - f(2)$
    D． $0 < f(3) - f(2) < f^{\prime}(2) < f^{\prime}(3)$
\end{quote}

8．已知函数 $y = f(x)$ 的图象如图所示，则 $f\prime\left( x_{A} \right)$ 与 $f\prime\left( x_{B} \right)$ 的大小关系是（~~~~）

\begin{quote}
    A． $f\prime\left( x_{B} \right) < f\prime\left( x_{A} \right) < 0$
    B． $f\prime\left( x_{A} \right) < f\prime\left( x_{B} \right) < 0$

    C． $f\prime\left( x_{A} \right) = f\prime\left( x_{B} \right)$
    D． $f\prime\left( x_{A} \right) > f\prime\left( x_{B} \right) > 0$
\end{quote}

9．若函数 $f(x) = 6x^{3} - 2f^{\prime}(1)x$ ，则 $f^{\prime}(1) =$ （~~~~）

\begin{quote}
    A．3 B．6 C． $- 6$  D． $- 3$
\end{quote}

10．已知曲线 $y = f(x)$ 在 $x = 5$ 处的切线方程是 $y = - 2x + 8$ ，则 $f(5) + f^{\prime}(5)$ 的值为（~~~）

\begin{quote}
    A． $- 4$  B． $- 2$  C． $2$  D． $4$
\end{quote}

11．已知点 $P$ 是曲线 $y = x^{2} - 3\text{ln}x$ 上任一点，则 $P$ 到直线 $x + y + \frac{3}{2} = 0$ 的距离的最小值是（~~~~）

\begin{quote}
    A． $\frac{7}{4}$  B． $\frac{7\sqrt{2}}{4}$  C． $\frac{3\sqrt{2}}{2}$
    D． $\frac{7}{8}$
\end{quote}

12．若函数 $f(x) = \frac{1}{\tan x}$ ，则导函数 $f^{\prime}(x) =$ （~~~）

\begin{quote}
    A． $\frac{1}{\sin^{2}x}$  B． $\frac{1}{\cos^{2}x}$
    C． $- \frac{1}{\sin^{2}x}$  D． $- \frac{1}{\cos^{2}x}$
\end{quote}

13．下列函数的求导不正确的是（~~~~）

\begin{quote}
    A． $\lbrack{(1 - 2x)}^{3}\rbrack^{\prime} = - 6{(1 - 2x)}^{2}$
    B． $\text{(cos}\frac{x}{3})^{\prime} = - \frac{1}{3}\text{sin}\frac{x}{3}$

    C． $(2^{2x + 1})^{\prime} = 2^{2x + 1}\text{ln}2$
    D． $\text{(e}^{2x})^{\prime} = 2\text{e}^{2x}$
\end{quote}

14．已知函数 $f(x) = 1 - 2x - \sin x$ ，则曲线 $y = f(x)$ 在点 $\left( 0,f(0) \right)$ 处的切线方程为（~~~~）

\begin{quote}
    A． $3x + y - 1 = 0$  B． $2x - y + 1 = 0$  C． $3x - y + 1 = 0$
    D． $2x + y - 1 = 0$
\end{quote}

15．函数 $y = - \frac{\cos x}{x}$ 的图像可能是（~~~~）

\begin{quote}
\end{quote}

16．函数 $f(x) = \frac{x^{3}}{\text{e}^{x} + 1}$ 的图象大致是（~~~~）

\begin{quote}
\end{quote}

17．函数 $f(x) = x^{3}\text{e}^{x}$ 的图象大致为

\begin{quote}
\end{quote}

18．已知向量 $\overrightarrow{a} = \left( - 1,\text{ }\sqrt{3} \right),\text{ }\left| \overrightarrow{b} \right| = 2\sqrt{3},\text{ }\overrightarrow{a} \cdot \left( \overrightarrow{a} - 2\overrightarrow{b} \right) = 16$ ，则向量 $\overrightarrow{a}$ 与 $\overrightarrow{b}$ 的夹角为（~~~~）

\begin{quote}
    A． $\frac{\text{π}}{3}$  B． $\frac{\text{π}}{6}$
    C． $\frac{5\text{π}}{6}$  D． $\frac{2\text{π}}{3}$
\end{quote}

\textbf{二、多选题}

19．下列函数求导正确的是（~~~）

\begin{quote}
    A．已知 $f(x) = x\ln x + x$ ，则 $f^{\prime}(x) = 2 + \ln x$

    B．已知 $f(x) = \text{e}^{2x}$ ，则 $f^{\prime}(x) = 2\text{e}^{2x}$

    C．已知 $f(x) = \frac{1}{x}$ ，则 $f^{\prime}(x) = \ln x$

    D．已知 $f(x) = x^{3} + \sin x$ ，则 $f^{\prime}(x) = 3x^{2} + \cos x$
\end{quote}

20．下列函数的导数计算正确的是（~~~~）

\begin{quote}
    A． $f(x) = xe^{x}，f^{'}(x) = (x + 1)e^{x}$
    B． $f(x) = \sqrt{x} + 1，f^{'}(x) = \frac{\sqrt{x}}{x}$

    C． $f(x) = \sin 2x$ ， $f^{'}(x) = 2\cos 2x$
    D． $f(x) = (2x + 1)^{5}，f^{'}(x) = 10(2x + 1)^{4}$
\end{quote}

\textbf{三、填空题}

21．已知 $f(x)$ 为偶函数，曲线 $y = f(x)$ 在点 $\left( 2,f(2) \right)$ 处的切线与直线 $4x - y - 3 = 0$ 垂直，设 $f(x)$ 的导函数为 $f^{\prime}(x)$ ，则 $f^{\prime}( - 2) =$
.

22．若向量 $\overrightarrow{a} = ( - 3 - 2t,2 - t)$ 与 $\overrightarrow{b} = (1, - 1)$ 的夹角为钝角，则实数 $t$ 的取值范围为
.

\textbf{四、解答题}

23．已知函数 $f(x) = x^{2}$ .

(1)求曲线 $y = f(x)$ 在点 $(2,f(2))$ 处的切线方程；

(2)求曲线 $y = f(x)$ 过点 $(2,0)$ 的切线方程.

24．已知函数 $f(x) = x^{3} + x - 16$ ．

(1)直线 $l$ 为曲线 $y = f(x)$ 在点 $(2, - 6)$ 处的切线，求直线 $l$ 的方程；

(2)求过原点且与曲线 $y = f(x)$ 相切的切线方程及切点坐标．

25．已知函数 $f(x) = \text{e}^{x - 1} - 2x - 1$ ，其导函数为 $f^{'}(x)$ .

(1)求 $f(x)$ 的单调区间；

(2)证明： $xf^{\prime}(x) + x - \ln x \geq 0$ .

26．已知函数  $f(x) = ax\text{ln}x - 2x + a,a \in \text{R}$  .

(1)若  $a = 1$  ，求曲线  $y = f(x)$  在点 $(1, - 1)$ 处的切线方程；

(2)若  $a > 0$  ，  $f(x) \geq 0$  恒成立，求  $a$  的取值范围.

27．已知单位向量 ${\overrightarrow{e}}_{1}$ ， ${\overrightarrow{e}}_{2}$ ，夹角为 $\frac{\text{π}}{2}$ ， $\overrightarrow{a} = - {\overrightarrow{e}}_{1}$ ， $\overrightarrow{b} = m{\overrightarrow{e}}_{1} + {\overrightarrow{e}}_{2}$ ，且 $\overrightarrow{a}$ 与 $\overrightarrow{b}$ 的夹角为 $\frac{\text{π}}{4}$ ．

(1)若 $\overrightarrow{a} + \lambda\overrightarrow{b}$ 与 $2\overrightarrow{a} + \overrightarrow{b}$ 所成的角为锐角，求实数 $\lambda$ 的取值范围；

(2)若向量 $\overrightarrow{c}$ 为 $\overrightarrow{a}$ 在 $\overrightarrow{b}$ 上的投影向量，求 $\left| \overrightarrow{a} + \overrightarrow{c} \right|$ ．
